\documentclass[12pt,a4paper]{article}
\usepackage[T1]{fontenc}
\usepackage{mathptmx}
\usepackage[margin=1in]{geometry}
\usepackage{amsmath}
\title{AI-evolved biological factories: constraint-based evolutionary search for robust metabolic architectures in \emph{E. coli}}
\author{}
\date{September 26, 2026}
\begin{document}
\maketitle

\section{Introduction}
Engineered microbes convert renewable feedstocks into fuels, materials, and chemicals, but the design of production strains remains guided primarily by human intuition: an engineer picks a pathway, expresses it, and hopes the host tolerates it. An alternative is to let a search process explore the space of possible metabolic architectures and select for what works, under conditions that resemble an industrial process rather than a single idealized optimum. This project asks whether a computational evolutionary search, starting from a target product rather than an existing pathway, can discover non-obvious biological architectures that match or exceed known engineered pathways under multi-condition robustness criteria. The design principles that let engineered cells maintain high productivity while remaining metabolically stable are the scientific object of interest; the evolved architecture library is the deliverable.

Three hypotheses are pre-registered (PRE-REGISTRATION.md, locked before any data or search): H1 (performance) - at least one evolved architecture reaches predicted product flux at or above the published engineered pathway evaluated in the same genome-scale model under identical uptake bounds; H2 (novelty) - the evolved library contains architectures absent from the published pathway that are mechanistically interpretable; H3 (robustness principle) - architectures selected under multi-condition pressure retain their production across a locked stress battery at a significantly higher rate than architectures selected under single-condition pressure. Honest negatives are preserved and redirect the project internally; no negative is published as the final result.

\section{Pre-registered design}
\subsection{Hosts, models, and targets}
The primary host is \emph{Escherichia coli} K-12 MG1655 in the genome-scale metabolic model iML1515, with the yeast arm (Yeast8) reserved for cross-host robustness analysis. Three locked target products span distinct chemistry classes: 1,4-butanediol (non-natural industrial diol; benchmark route Yim et al.\ 2011, the Genomatica pathway), isobutanol (branched-chain higher alcohol; benchmark Atsumi et al.\ 2008, the keto-acid route), and lycopene (tetraterpene; benchmark Alper et al.\ 2005, crtEBI with dxs push). Reaction lists per the cited papers were locked in data/benchmarks.json before any search, with DOI provenance.

\subsection{Benchmark metric: production envelopes}
For each target, the published engineered pathway was implemented in iML1515 under the locked base condition (M9 minimal, glucose 10 mmol/gDW/h, aerobic). Because production competes with biomass, a single max-growth number is meaningless: at 100\% of maximum growth all three pathways carry zero product flux. The benchmark is therefore a production envelope: maximum product flux with growth fixed at locked fractions of the FBA maximum,
\begin{equation}
F_t(g) = \max_{v} \; v_{\text{product}} \quad \text{s.t.} \quad S v = 0, \; v_{\text{biomass}} = g \cdot \mu_{\max}, \; l \le v \le u,
\end{equation}
evaluated at $g \in \{1.0, 0.9, 0.75, 0.5, 0.25, 0.0\}$ with a parsimonious-FBA tie-break. Measured envelopes (results/benchmark\_flux.json): 1,4-BDO reaches $F=10.35$ mmol/gDW/h at $g=0$ and $5.27$ at $g=0.5$; isobutanol $10.00$ and $5.24$; lycopene $0.97$ and $0.49$. The envelope, not a single titer, is what an evolved architecture must beat at matched growth.
Table~\ref{tab:envelopes} reports the measured envelopes for all three published routes at the six locked growth fractions; the FBA maximum growth rate of the host under the base condition is $\mu_{\max} = 0.877$ h$^{-1}$ for all targets (identical uptake bounds, wild-type core).
\begin{table}[h]
\centering
\caption{Measured production envelopes (product flux, mmol/gDW/h) of the three published engineered routes in iML1515 under the locked base condition, at growth fixed to fractions of $\mu_{\max}$ (results/benchmark\_flux.json, run 2026-09-26; pFBA tie-break coefficient 0.999).}
\label{tab:envelopes}
\begin{tabular}{lcccccc}
\hline
Target route & $g{=}1.0$ & $0.9$ & $0.75$ & $0.5$ & $0.25$ & $0.0$ \\
\hline
1,4-BDO (Yim 2011) & 0.000 & 1.060 & 2.649 & 5.267 & 7.878 & 10.355 \\
Isobutanol (Atsumi 2008) & 0.000 & 1.048 & 2.620 & 5.239 & 7.813 & 10.000 \\
Lycopene (Alper 2005) & 0.000 & 0.098 & 0.244 & 0.488 & 0.731 & 0.970 \\
\hline
\end{tabular}
\end{table}
The zero at $g=1.0$ in every row is a verified property of the host--route system, not a solver artifact: biomass and product draw on the same carbon and redox budget, so the envelopes are strictly monotone decreasing in $g$, and any comparison at a single growth point would be meaningless without the matched-growth rule.

\subsection{Search protocol}
The search genome encodes on/off choices over a locked reaction pool (benchmark route blocks plus literature-noted alternatives, data/reaction\_pool.json). Fitness is product flux at the pFBA-optimal solution with growth fixed at $g=0.5$, the envelope midpoint. The search is a ($\mu+\lambda$) evolutionary algorithm ($\mu=24$, $\lambda=48$, tournament selection, 60 generations, five locked seeds), with a 1000-seed random-architecture null for significance calibration. On the version-1 pool (two alternative blocks per target, four genomes), exhaustive enumeration strictly dominates the stochastic search and was used instead; the EA scaffold is retained for expanded pools.

\subsection{Stability battery}
Robustness is evaluated over a locked nine-condition battery (nutrient shifts, carbon shifts to glycerol and acetate, anaerobiosis, uptake perturbations, oxygen limitation). After a metric-interpretation bug was caught at first evaluation (growth referenced to the base-condition maximum, which made every condition untestable), the locked rule is growth fixed at 50\% of the condition-specific maximum, with retention defined against nominal-condition flux. The nine locked conditions are: GLC\_LOW (glucose uptake $-2$), GLC\_MID ($-5$), GLC$+20\%$ ($-12$), GLC$-20\%$ ($-8$), GLYCEROL (sole carbon, uptake $-10$), ACETATE (sole carbon, uptake $-10$), ANAEROBIC (oxygen uptake 0), O2$-50\%$ (oxygen uptake halved), and the nominal base GLC\_AEROBIC (glucose $-10$, aerobic); all uptakes in mmol/gDW/h. Formally, for architecture $a$ and condition $c$,
\begin{equation}
R(a, c) \;=\; \frac{F_a^{(c)}(0.5)}{F_a^{(\mathrm{nominal})}(0.5)}, \qquad
\mathrm{BatteryPass}(a) \;=\; \bigwedge_{c \in \mathrm{battery}} \mathbf{1}\!\left[ R(a,c) \ge 0.8 \right],
\end{equation}
where $F_a^{(c)}(0.5)$ is the architecture's product flux in condition $c$ with growth fixed at $50\%$ of the condition-specific maximum. The pass bar of $0.8$ was locked before any architecture was evaluated against the battery.

\section{Results}
\subsection{Single-condition search reproduces the benchmark envelopes}
Across all fifteen seed runs (three targets, five seeds), the best evolved architecture exactly reaches the locked envelope value at $g=0.5$: 5.27 mmol/gDW/h for 1,4-BDO, 5.24 for isobutanol, 0.49 for lycopene. The search recovers the published routes as optimal under single-condition selection - a validation of the evaluation pipeline, and the baseline for H1.

\subsection{H2 candidate: a tail-only minimal 1,4-BDO architecture}
A non-obvious architecture emerged from the search: with \emph{both} heterologous upstream route blocks switched off, the four-step heterologous tail (4-hydroxybutyrate dehydrogenase, CoA transferase, CoA-acylating aldehyde dehydrogenase, alcohol dehydrogenase) still reaches the full nominal flux of 5.27 mmol/gDW/h. The model explains why: \emph{E. coli}'s native GABA shunt (succinate-semialdehyde supply through glutamate metabolism) feeds the tail without any heterologous upstream enzyme. This tail-only architecture is absent from the Yim et al.\ pathway definition, uses four fewer heterologous enzymes, and is mechanistically interpretable - the H2 novelty pattern in its locked form, pending verification against the full battery and the random-architecture null.

\subsection{H3 honest negative on pool v1, and a battery-design finding}
Exhaustive multi-condition evaluation of pool v1 found that \emph{no} genome passes the locked battery (80\% retention in all nine conditions): best minimum retentions are 0.175 (1,4-BDO), 0.176 (isobutanol), and 0.27 (lycopene). Inspection shows the failure is structural, not a lack of search power: under the severe nutrient shift GLC\_LOW (glucose at 2 rather than 10 mmol/gDW/h), elementary carbon balance caps a C4 product at roughly one to two mmol/gDW/h against a nominal 5.27, a retention ceiling of 0.2--0.3 that no architecture can lift. The bound is elementary: a C4 product made from glucose at uptake $u_{\mathrm{glc}}$ satisfies
\begin{equation}
v_{\mathrm{product}} \;\le\; \frac{6}{4}\, u_{\mathrm{glc}} \;=\; 3.0 \quad \text{mmol/gDW/h at } u_{\mathrm{glc}} = 2,
\end{equation}
before any carbon is allocated to biomass, maintenance, or byproducts; with growth fixed at $50\%$ of the condition maximum the achievable product flux falls to roughly one to two, giving $R(a, \mathrm{GLC\_LOW}) \le 1.5/5.267 \approx 0.28$ for every conceivable architecture $a$ - including the published routes and any architecture the search could ever return. The benchmark engineered routes fail identically. A pass bar of 80\% retention on severe nutrient shifts is therefore unsatisfiable by construction, the battery provides no selection gradient, and H3 is untestable as locked. This is recorded as an honest negative with the design flaw named, and a redesign (retention against condition-achievable optimum rather than against nominal flux; tiered bars for shifts versus perturbations; robustness-oriented pool levers such as knockout blocks and co-utilization feeds) is being locked as a dated amendment before any re-evaluation, following the project's redirect-on-failure rule.

\subsection{Pool v2: robustness levers shift the nominal optimum (amendment 2026-09-26)}
Following the redirect rule, pool v2 was locked as a dated amendment before evaluation: four host-level robustness levers - acetate-overflow knockout (ackA--pta), formate-overflow knockout (pflB), lactate-overflow knockout (ldhA), and a glycerol co-utilization feed - added to the target-specific blocks, expanding the genome space to 64 genomes per target. Exhaustive enumeration (192 genomes total, metric-independent of the pending battery redesign) was again strictly dominant over the stochastic budget at this size. Under the nominal condition at $g=0.5$, the co-utilization feed dominates every target: best nominal fluxes rise to 6.496 (1,4-BDO), 6.498 (isobutanol), and 0.730 (lycopene) mmol/gDW/h, against envelope values of 5.267, 5.239, and 0.488 - nominal-condition exceedances of 23\%, 24\%, and 50\% over the published routes evaluated in the same model. The gain is attributable to the added carbon supply, so the honest comparison under co-feeding must re-baseline the benchmark route with the same feed block enabled; that re-baselined comparison is part of the pending battery-redesign amendment. An honest secondary negative is recorded: for isobutanol, some knockout combinations abolish product flux entirely (minimum 0.0), so the levers are not uniformly benign even at nominal condition.

\section{Caveats}
(1) All results are in-silico flux predictions; no wet-lab validation is claimed, and no build-phase engineering guidance is produced. (2) Both pools evaluated so far are deliberately small; the exhaustive results bound what these pools can express, and conclusions about search power await larger pools. (3) The tail-only H2 candidate is verified at nominal condition only; its robustness profile is part of the battery redesign. (4) GSMM predictions assume steady state and fixed biomass composition; real titers are typically lower.

\section{Next steps}
(1) Lock the battery-redesign amendment from the staged redirection consult and re-evaluate H3. (2) Re-baseline the benchmark routes with the co-feed block enabled and re-run the full EA at the locked budget on pool v2, with the 1000-seed random null. (3) Expand the pool further with cofactor and transporter variants, each with provenance. (4) Yeast8 cross-host arm for the H3 robustness principle.

\end{document}
